\documentclass[12pt]{article}
\usepackage[utf8]{inputenc}
\usepackage[T1]{fontenc}
\usepackage[margin=1in]{geometry}
\usepackage{setspace}
\usepackage{xcolor}
\usepackage[colorlinks=true, urlcolor=blue!60!black, linkcolor=blue!60!black]{hyperref}
\setlength{\parindent}{2em}
\setlength{\parskip}{1.25em}
\renewcommand{\baselinestretch}{1.0}
\begin{document}
\begin{center}
{\LARGE \textit{City of Boerne v. Flores}, 521 U.S. 507 (1997)\par}
\vspace{0.8em}
{\large Case Brief --- Four-Page Edition\par}
\end{center}
\vspace{1em}

\section{Memory Jogger}

The Court struck down the Religious Freedom Restoration Act of 1993 as applied to the states, holding that Congress's Section 5 enforcement power is remedial and preventive---requiring ``congruence and proportionality'' between injury and means---rather than a power to redefine constitutional rights (\href{https://supreme.justia.com/cases/federal/us/521/507/}{521 U.S. 507}).

\section{Detailed Case Facts}

St. Peter Catholic Church in Boerne, Texas, seated about 230 worshippers, too few for its growing parish. In 1993 Archbishop Patrick F. Flores authorized an expansion, but the city soon enacted a historic preservation ordinance and denied the building permit because the church lay within a newly designated historic district.

Flores sued under RFRA, which Congress had enacted by overwhelming bipartisan majorities in response to \textit{Employment Division v. Smith} (1990), where the Court held that the Free Exercise Clause does not exempt religious objectors from neutral, generally applicable laws. Aiming to ``restore the compelling interest test'' of \textit{Sherbert} and \textit{Yoder}, RFRA barred all levels of government from substantially burdening religious exercise unless the burden furthered a ``compelling governmental interest'' by the ``least restrictive means.'' Congress relied on Section 5 of the Fourteenth Amendment to apply RFRA to the states.

\section{Procedural History}

The Western District of Texas held that RFRA exceeded Congress's Section 5 power and violated separation of powers. The Fifth Circuit reversed, upholding RFRA as broader statutory protection above \textit{Smith}'s constitutional floor. The Supreme Court granted certiorari, heard argument February 19, 1997, and reversed on June 25, 1997.

\section{Judicial Votes}

Reversed, \textbf{6--3}.
\begin{itemize}
\item \textbf{Majority:} Justice Kennedy, joined by Chief Justice Rehnquist and Justices Stevens, Thomas, and Ginsburg, and by Justice Scalia except Part III-A-1 (drafting history).
\item \textbf{Concurrences:} Justice Stevens; Justice Scalia (in part), joined by Justice Stevens.
\item \textbf{Dissents:} Justice O'Connor (joined in part by Justice Breyer); Justice Souter; Justice Breyer.
\end{itemize}

\section{Holding}

RFRA exceeds Congress's Section 5 power and is unconstitutional as applied to state and local governments. Section 5's power ``to enforce'' is remedial and preventive, not substantive: Congress may not alter the meaning of constitutional rights as interpreted by the Court, and remedial legislation must show congruence and proportionality between the injury and the means. RFRA fails because its record lacks evidence of modern religious bigotry and its sweeping, perpetual strict-scrutiny mandate is out of proportion to any pattern of unconstitutional conduct. The Court did not invalidate RFRA as applied to the federal government.

\section{Analysis of Opinions}

\textbf{Kennedy (majority).} Section 5 is a ``positive grant of legislative power'' that allows preventive rules, but it grants only the power ``to enforce.'' The 39th Congress rejected John Bingham's first draft, which would have given Congress plenary power, in favor of self-executing limits and a corrective enforcement power. Kennedy rejected the ``ratchet theory'' reading of \textit{Katzenbach v. Morgan}, which would place the Constitution ``on a level with ordinary legislative acts.'' Unlike the Voting Rights Act, which answered a century of documented discrimination with targeted coverage, RFRA rested on anecdotes of incidental burdens, reached every law and agency without sunset or geographic limit, and imposed ``the most demanding test known to constitutional law.''

\textbf{Stevens (concurring).} RFRA violates the Establishment Clause by giving religious institutions ``a legal weapon that no atheist or agnostic can obtain.''

\textbf{Scalia (concurring in part).} Founding-era sources protected only against laws targeting religion; early accommodations show legislatures \textit{may} grant exemptions, not that courts must. The question is whether legislatures or courts decide.

\textbf{O'Connor (dissenting).} She accepted the Section 5 analysis but would have reargued and overruled \textit{Smith}, which she called ``demonstrably wrong,'' reading the Free Exercise Clause as an affirmative guarantee.

\textbf{Souter and Breyer (dissenting).} Souter had ``serious doubts'' about \textit{Smith} and would have ordered reargument or dismissed the writ; Breyer agreed on reargument without endorsing the Section 5 analysis.

\section{Examples: Future Applications}

\begin{enumerate}
\item \textbf{Same-Side (Invalidated): ADA Damages Against States.} A statute lets individuals sue states for failing to accommodate disabilities. Disability gets only rational-basis review, and without a record of unconstitutional state discrimination the remedy is not congruent and proportional (as in \textit{Garrett}, 2001).
\item \textbf{Same-Side (Invalidated): Redefining Takings.} A Section 5 statute requires compensation for any regulation reducing property value by 15\%. It rewrites the \textit{Penn Central} standard rather than remedying unconstitutional takings.
\item \textbf{Opposite-Side (Upheld): Family Leave.} A narrow leave mandate answers a documented pattern of gender discrimination, which triggers heightened scrutiny, so it is congruent and proportional (as in \textit{Hibbs}, 2003).
\item \textbf{Opposite-Side (Upheld): RLUIPA.} A narrower statute imposes strict scrutiny on prisons and zoning boards only when federal funds or interstate commerce are involved. Resting on Spending and Commerce powers, it avoids \textit{Boerne}'s bar (see \textit{Cutter v. Wilkinson}, 2005).
\item \textbf{Fence-Sitter: Nationwide Voting Rules.} A statute bars voter-ID rules producing racial disparities nationwide, without sunset or geographic limits. Race and voting earn Congress latitude under \textit{South Carolina v. Katzenbach}, but the missing tailoring invites a serious challenge.
\end{enumerate}

\section{Critique}

\textbf{Progressive and originalist critics.} Robert Post and Reva Siegel fault \textit{Boerne}'s ``juricentric'' model, which treats Congress like an inferior agency building a trial record, and note that the test enabled \textit{Morrison}, \textit{Kimel}, and \textit{Garrett}. Michael McConnell argues that ``appropriate legislation'' signaled deferential \textit{McCulloch} review, that the Reconstruction Congress distrusted the Court after \textit{Dred Scott}, and that Kennedy misread the Bingham draft's rejection.

\textbf{Federalism defenders.} Jeffrey Sutton and William Van Alstyne praise \textit{Boerne} as a guardrail: otherwise Congress would gain a general police power and ordinary majorities could amend the Constitution without Article V, and RLUIPA shows Congress can still act through enumerated Article I powers.

\section{Key Quotations}

\begin{quote}
``There must be a congruence and proportionality between the injury to be prevented or remedied and the means adopted to that end. Lacking such a connection, legislation may become substantive in operation and effect.''
\end{quote}
\noindent --- Justice Kennedy, Opinion of the Court, 521 U.S. at 520.

\begin{quote}
``Legislation which alters the meaning of the Free Exercise Clause cannot be said to be enforcing the Clause. Congress does not enforce a constitutional right by changing what the right is. It has been given the power `to enforce,' not the power to determine what constitutes a constitutional violation.''
\end{quote}
\noindent --- Justice Kennedy, Opinion of the Court, 521 U.S. at 519.

\begin{quote}
``Shifting legislative majorities could change the Constitution and effectively circumvent the difficult and detailed amendment process contained in Article V.''
\end{quote}
\noindent --- Justice Kennedy, Opinion of the Court, 521 U.S. at 529.

\end{document}
